% ===================================================================
%  TAULA N.º 10 — El Mar. — El Port.
%  Traduction 2026 d'apres texto/io, controlee sur texto/fr et
%  texto/es. Consigne : voir texto/ca/10-taula-01.tex.
%
%  Le francais decale sa numerotation d'un cran a partir de l'alinea
%  8 ; c'est une coquille de l'atelier francais, et la colonne suit
%  l'ido, qui compte juste.
% ===================================================================

%%K t10-apar-1 apar
\VUsaut{4.0mm}
\VUtitre{15.0pt}{60}{TAULA N.º 10}
\VUsaut{3.0mm}
\VUpk{11.4pt}{40}{El Mar. --- El Port.}
\VUsaut{3.0mm}

%%K t10-01-1 p
1. --- A l'estiu, mentre uns van a buscar la calma i la fresca a la muntanya,
d'altres prefereixen anar a la vora del mar. Això és el que vam fer l'any passat,
perquè ens agrada molt l'\VUgras{oceà} \textsuperscript{(1)}.

%%K t10-02-1 p
2. --- Per la meva banda, sento un gran plaer en contemplar aquella immensa capa
d'aigua que s'estén fins a l'\VUgras{horitzó} \textsuperscript{(2)} i en veure
partir per a una llarga \VUgras{travessia} els enormes \VUgras{vaixells de vapor}
\textsuperscript{(3)} en què embarquen els viatgers que van a Amèrica.

%%K t10-03-1 p
3. --- Per això el meu pare, que coneix els meus gustos, tria sempre, per passar
les vacances, una platja molt tranquil·la situada vora un dels nostres grans
\VUgras{ports de guerra}.

%%K t10-03-2 p
Durant la nostra darrera estada, l'\VUgras{esquadra} havia vingut a llançar
l'àncora darrere de l'enorme \VUgras{dic} \textsuperscript{(4)} que tanca i
protegeix el \VUgras{port militar} \textsuperscript{(5)}, i vaig poder admirar al
meu gust els diferents \VUgras{vaixells de guerra}, des del petit
\VUgras{submarí} \textsuperscript{(10)}, que se submergeix i desapareix sota les
aigües, fins a l'enorme \VUgras{cuirassat} \textsuperscript{(9)}, que fins ara
gairebé només temia els atacs del \VUgras{torpediner} \textsuperscript{(8)}.

%%K t10-tit-2 sub
\VUsaut{3.0mm}
\VUpk{10.2pt}{40}{Les embarcacions petites.}
\VUsaut{3.0mm}

%%K t10-04-1 p
4. --- Aquest ja sabia descobrir molt \VUgras{hàbilment} el punt feble de la
gruixuda cuirassa metàl·lica per clavar-hi el perillós \VUgras{torpede}, l'explosió
del qual causa la pèrdua del vaixell; però era menys temible que el submarí,
perquè els contratorpediners li donaven caça fàcilment.

%%K t10-05-1 p
5. --- També vaig poder veure els ràpids \VUgras{creuers} \textsuperscript{(6)}
cuirassats, gairebé tan invulnerables com el veritable cuirassat, diversos
\VUgras{transports} \textsuperscript{(12)}, tres o quatre \VUgras{canoners}
\textsuperscript{(7)} i, finalment, una \VUgras{fragata} \textsuperscript{(11)}.
\VUgras{L'espectacle d'aquella flota, quan maniobra per llevar àncores, és dels
més interessants.}

%%K t10-06-1 p
6. --- Quina diferència de construcció entre aquelles enormes masses d'acer i els
elegants \VUgras{bergantins de tres pals} o els fins \VUgras{velers}
\textsuperscript{(13)}, que encara s'usen a la marina mercant, i que jo veia
entrar al port a tota vela quan anava a passejar pel \VUgras{moll}
\textsuperscript{(14)}! És cert que aquests perdien molt dels seus avantatges quan
el \VUgras{vent} era contrari. Havien de recollir totes les veles per entrar a la
dàrsena, i calia un \VUgras{remolcador} \textsuperscript{(16)} que anés a
buscar-los per portar-los al moll com a simples \VUgras{gavarres}
\textsuperscript{(17)}.

%%K t10-tit-3 sub
\VUsaut{3.0mm}
\VUpk{10.2pt}{40}{El port comercial.}
\VUsaut{3.0mm}

%%K t10-07-1 p
7. --- L'interessant del port de què parlo és que no comprèn només un port de
guerra, creat artificialment amb obres gegantines, sinó a més un meravellós
\VUgras{port de comerç} situat a la \VUgras{desembocadura} d'un gran riu.

%%K t10-08-1 p
8. --- La llengua de terra que separa les dues dàrsenes està fortificada i
defensada per una sèrie de \VUgras{bateries} i \VUgras{baluards} que
s'endinsen al mar. Cal un permís especial per passejar per les \VUgras{muralles}
\textsuperscript{(15)}. Jo l'havia obtingut sense dificultat per mediació del meu
oncle, que és enginyer de les \VUgras{drassanes} \textsuperscript{(18)}, i vaig
anar a visitar els vaixells que es construeixen sota vasts coberts.

%%K t10-09-1 p
9. --- Fins i tot vaig assistir a la botadura d'un cuirassat i recordo el moment
d'emoció que va sentir tota la multitud quan l'enorme massa, abandonada a si
mateixa, va començar a lliscar per la seva grada i va venir a surar damunt
l'aigua del riu.

%%K t10-10-1 p
10. --- Per anar a les drassanes es creua el \VUgras{camp de maniobres}
\textsuperscript{(19)}, on les tropes de la guarnició vénen cada dia a fer
l'exercici, i es passa per una \VUgras{passarel·la} \textsuperscript{(20)} de ferro
que uneix l'esplanada amb l'\VUgras{arsenal} \textsuperscript{(21)}.

%%K t10-tit-4 sub
\VUsaut{3.0mm}
\VUpk{10.2pt}{40}{L'Arsenal i els dics.}
\VUsaut{3.0mm}

%%K t10-11-1 p
11. --- Amb el meu oncle vaig visitar aquell gran establiment ple de
\VUgras{canons} \textsuperscript{(22)}, \VUgras{bales} \textsuperscript{(23)} i
\VUgras{obusos}. Vaig veure també la \VUgras{fàbrica metal·lúrgica}
\textsuperscript{(24)} en què es fabriquen les \VUgras{calderes}
\textsuperscript{(25)} dels vaixells de vapor.

%%K t10-12-1 p
12. --- Molt a prop hi ha els \VUgras{dics} \textsuperscript{(26)} --- dics de
carena --- i els curiosíssims \VUgras{dics flotants} \textsuperscript{(27)}, en
què vaig poder veure de prop, en un vaixell en reparació, l'\VUgras{hèlix}
\textsuperscript{(28)}, la \VUgras{quilla} \textsuperscript{(29)} i el
\VUgras{timó} \textsuperscript{(30)} d'una nau.

%%K t10-13-1 p
13. --- El port de comerç és tan interessant d'estudiar com el port de guerra, i
vaig passar moltes hores badant pels molls on són amarrats els \VUgras{vapors de
rodes} \textsuperscript{(31)} que remunten el riu, i mirant funcionar la
\VUgras{grua d'arbrar} \textsuperscript{(32)}, que aixeca com plomes càrregues
enormes.

%%K t10-tit-5 sub
\VUsaut{4.0mm}
\VUpk{10.2pt}{40}{El Pràctic.}
\VUsaut{3.0mm}

%%K t10-14-1 p
14. --- Admirava els \VUgras{transatlàntics} \textsuperscript{(34)} que sortien de
la rada amb les seves \VUgras{banderes} \textsuperscript{(35)} al vent, conduïts
per un hàbil \VUgras{pràctic} \textsuperscript{(36)} que sabia meravellosament
seguir el \VUgras{canal} marcat per \VUgras{boies} \textsuperscript{(40)} i
\VUgras{balises}, evitant les \VUgras{gavarres} \textsuperscript{(41)} que es
dirigeixen amb tanta dificultat amb la seva única vela quadrada. Distingia a la
\VUgras{proa} \textsuperscript{(37)} els \VUgras{camarots}
\textsuperscript{(39)} de segona classe i, a la \VUgras{popa}
\textsuperscript{(38)}, els salons dels passatgers de primera.

%%K t10-15-1 p
15. --- De vegades assistia també a la càrrega d'un \VUgras{vaixell de càrrega}
\textsuperscript{(42)}, en què s'amuntegaven les mercaderies del
\VUgras{magatzem} \textsuperscript{(43)} i de l'\VUgras{estació marítima}
\textsuperscript{(44)}, portades pels nombrosos trens de la \VUgras{línia dels
molls} \textsuperscript{(45)}.

%%K t10-tit-6 sub
\VUsaut{3.0mm}
\VUpk{10.2pt}{40}{Remar per gust.}
\VUsaut{3.0mm}

%%K t10-16-1 p
16. --- Sovint remava a la \VUgras{dàrsena a flot} \textsuperscript{(53)}, on es
refugien les petites \VUgras{embarcacions} \textsuperscript{(46)}, i era una
alegria per a mi de manejar el \VUgras{rem} \textsuperscript{(47)}. Pujava a la
\VUgras{coberta} \textsuperscript{(48)} d'una \VUgras{barca de vela}
\textsuperscript{(49)}, enfilava amb ajuda del \VUgras{cordam}
\textsuperscript{(52)} pel \VUgras{pal} \textsuperscript{(51)} fins a les
\VUgras{vergues} \textsuperscript{(50)}, i després reprenia el meu passeig en
\VUgras{iol} \textsuperscript{(54)} entre les pesades \VUgras{barques de pesca}
\textsuperscript{(55)}, on s'assequen les \VUgras{xarxes}
\textsuperscript{(56)}.

%%K t10-17-1 p
17. --- Pel \VUgras{moll} \textsuperscript{(58)} desfilaven davant meu les
\VUgras{mercaderies} \textsuperscript{(59)} més diverses, ficades en
\VUgras{caixes} \textsuperscript{(60)} o en \VUgras{fardells}
\textsuperscript{(61)}. Formaven el \VUgras{carregament} \textsuperscript{(63)} de
les barcasses, i potents \VUgras{grues} \textsuperscript{(62)} les aixecaven i les
deixaven en \VUgras{camions} \textsuperscript{(64)} mentre els
\VUgras{transportistes} feien la declaració i pagaven els drets a la
\VUgras{duana} \textsuperscript{(65)}.

%%K t10-18-1 p
18. --- Finalment, quan arribava al \VUgras{pont giratori} \textsuperscript{(66)} o
a la \VUgras{resclosa} \textsuperscript{(69)}, passant per davant de la
\VUgras{draga} \textsuperscript{(68)} que neteja el \VUgras{canal}
\textsuperscript{(67)}, desembarcava al moll, vora el \VUgras{semàfor}
\textsuperscript{(70)}.

%%K t10-19-1 p
19. --- És a l'altre costat d'aquell moll on atraquen les \VUgras{xalupes}
\textsuperscript{(71)} que porten a terra els oficials de l'esquadra. És un
espectacle bell veure els mariners aixecar els seus rems a compàs mentre
l'embarcació, portada per l'\VUgras{ona} \textsuperscript{(72)}, aborda sense
esforç l'escala del desembarcador. És aquí on millor es pot observar la
\VUgras{marea} i el moviment del \VUgras{flux} i del \VUgras{reflux} a la
\VUgras{platja} \textsuperscript{(73)}.

%%K t10-tit-7 sub
\VUsaut{3.0mm}
\VUpk{10.2pt}{40}{El Casino d'oficials.}
\VUsaut{3.0mm}

%%K t10-20-1 p
20. --- Per tornar a casa agafava el \VUgras{tramvia elèctric}
\textsuperscript{(74)}, la \VUgras{parada} \textsuperscript{(75)} del qual és al
\VUgras{pal} que hi ha davant del casino dels oficials.

%%K t10-20-2 p
Des del \VUgras{balcó} \textsuperscript{(76)} del casino, adornat amb
\VUgras{àncores} \textsuperscript{(77)}, canons i bales, veia sovint una brillant
reunió d'oficials de marina: \VUgras{tinents de navili} \textsuperscript{(78)} amb
la seva espasa, \VUgras{capitans d'artilleria} \textsuperscript{(79)} i
d'\VUgras{infanteria} \textsuperscript{(80)} de marina amb el seu \VUgras{sabre}.
Darrere d'ells hi havia un \VUgras{mariner} \textsuperscript{(81)} que els portava
una \VUgras{ullera de llarga vista} quan volien distingir el que passava lluny.

%%K t10-21-1 p
21. --- Veia també joves \VUgras{alferes de navili} \textsuperscript{(82)} que
rebien les ordres del seu \VUgras{comandant} \textsuperscript{(83)} o de
l'\VUgras{almirall} \textsuperscript{(84)}, mentre un \VUgras{caporal}
\textsuperscript{(85)} esperava per portar-les a bord.

%%K t10-22-1 p
22. --- Des d'aquell balcó la vista és magnífica: es domina tota la platja amb les
seves \VUgras{tendes} \textsuperscript{(86)} i les seves \VUgras{casetes}
\textsuperscript{(87)}, i es veu, a l'hora del bany, els \VUgras{banyistes}
\textsuperscript{(88)} que juguen alegrement davant del \VUgras{casino}
\textsuperscript{(89)}.

%%K t10-23-1 p
23. --- A l'extrem de la platja, la costa forma una petita \VUgras{península}
\textsuperscript{(91)} rocosa, on ve a trencar l'\VUgras{onatge}
\textsuperscript{(92)} i damunt la qual està construït un \VUgras{far}
\textsuperscript{(93)} que indica de nit el camí als navegants. Aquella península
només està unida a terra per un \VUgras{istme} \textsuperscript{(90)} molt
estret.

%%K t10-tit-8 sub
\VUsaut{3.0mm}
\VUpk{10.2pt}{40}{Vista general de la costa.}
\VUsaut{3.0mm}

%%K t10-24-1 p
24. --- El \VUgras{cingle} \textsuperscript{(94)} es prolonga, retallat pel mar, que
hi dibuixa capriciosament una sèrie de \VUgras{badies} \textsuperscript{(95)} i
\VUgras{promontoris} (caps) que el fan dels més pintorescos.

%%K t10-24-2 p
A l'\VUgras{altiplà} \textsuperscript{(96)}, que domina l'\VUgras{estret}
\textsuperscript{(98)}, hi ha un \VUgras{fort} \textsuperscript{(97)} molt
important i alguns xalets.

%%K t10-25-1 p
25. --- Aquell estret separa del continent una \VUgras{illa}
\textsuperscript{(99)} molt perillosa, envoltada d'\VUgras{esculls}
\textsuperscript{(100)} i \VUgras{esculleres}, on barques i fins i tot grans
vaixells encallen de vegades. Malgrat la promptitud dels socors i de la ràpida
arribada dels \VUgras{bots salvavides}, aquells \VUgras{naufragis} són terribles
i, en les fortes \VUgras{tempestes} d'equinocci, diversos vaixells s'han perdut
amb homes i mercaderies.

%%K t10-25-2 p
Malgrat aquell veïnatge, \VUgras{el port} el freqüenten molt els mariners.
\VUsaut{6.0mm}
\VUfilet{22mm}
